\Needspace{8\baselineskip}\originalchapter{2006 作业七: 碰撞中心与波反射}
题源: 15 Nov 2006, due 22 Nov.
\begin{exercise}\label{ex:06-7-1}
原题1. 证明均匀B内一次离子-电子碰撞时两个导引中心的位移相等. 此断言在课程Lecture15给出.
\end{exercise}
\begin{solution}
碰撞短, 位置改变可忽略, 中心 $\mathbf r_g=\mathbf r+\mathbf p\times\mathbf B/(qB^2)$, 故 $\Delta r_g=\Delta p\times B/(qB^2)$. 动量守恒 $\Delta p_i=-\Delta p_e$, 单电荷 $q_i=-q_e$, 两负号抵消, 所以相等. q大小不同的多价离子时需按电荷权重写关系, 不是每个中心都一样.
\end{solution}
\begin{exercise}\label{ex:06-7-2}
原题2. 估横向热扩散. 同m、q粒子, 第1初V沿x, 第2初静, 碰撞中B忽略且rL\ensuremath{\gg}\ensuremath{\lambda}D. (a)第1出速度与x夹角\ensuremath{\theta}, 求两实验室速度; (b)碰撞位置位移小, B沿\ensuremath{-}z、q与\ensuremath{\omega}c定义正, 证明中心 $r_{g1}=\mathbf e_z\times\mathbf v_1/\omega_c$, rg2=0及碰后式; (c)按碰后能量加权中心定义\ensuremath{\Delta}rE; (d)以 $\chi\sim nV_{th}\int_0^\infty|\Delta r_E|^2\,2\pi b\dd b$ 求ln\ensuremath{\Lambda}阶, 哪个物种主导?
\end{exercise}
\begin{solution}
(a)动量和能量两式给 $\mathbf v_1'=V\cos\theta(\cos\theta\mathbf e_x+\sin\theta\mathbf e_y)$, $\mathbf v_2'=V(\sin^2\theta\mathbf e_x-\sin\theta\cos\theta\mathbf e_y)$; 能量分数cos\ensuremath{{}^{2}}、sin\ensuremath{{}^{2}}. (b)瞬间速度变而位置固定, 均匀回旋几何即所给中心, \ensuremath{\omega}c=q|B|/m.
(c)rL=V/\ensuremath{\omega}c, 权重代入得
\begin{align}
 \Delta r_{E,x}&=r_L\sin\theta\cos\theta(\sin^2\theta-\cos^2\theta),\\
 \Delta r_{E,y}&=-2r_L\sin^2\theta\cos^2\theta,\qquad
 |\Delta r_E|^2=r_L^2\sin^2\theta\cos^2\theta.
\end{align}
源提示x项写sin\ensuremath{{}^{4}}\ensuremath{-}cos\ensuremath{{}^{4}}, 因$\sin^2\theta+\cos^2\theta=1$与此同式. (d)同质量相对散射角2\ensuremath{\theta}, $\cot\theta=b/b_{90}$, $b_{90}=2K/(mV^2)$, K=q\ensuremath{{}^{2}}/(4\ensuremath{\pi}\ensuremath{\epsilon}0). 小角平方\ensuremath{\simeq}rL\ensuremath{{}^{2}}b90\ensuremath{{}^{2}}/b\ensuremath{{}^{2}}, 积分得
$\chi\sim2\pi nVr_L^2b_{90}^2\ln\Lambda=8\pi nK^2\ln\Lambda/(m^2\omega_c^2V)$.
同T用V\ensuremath{\propto}m$^{-1/2}$, 得$\chi_\perp\propto\sqrt m$, 所以离子大于电子约$\sqrt{m_i/m_e}$. 这是题定随机步长估计, 精确输运还含碰撞配对及速度平均系数.
\end{solution}
\begin{exercise}\label{ex:06-7-3}
原题3. 正入射未磁化等离子体界面, \ensuremath{\epsilon}=\ensuremath{\epsilon}0(1\ensuremath{-}\ensuremath{\omega}p\ensuremath{{}^{2}}/\ensuremath{\omega}\ensuremath{{}^{2}}). (a)求复电场反射系数\ensuremath{\Gamma}=Eref/Einc, 何时全反射, n=$10^{12}$m$^{-3}$数值; (b)沿z密度渐增、E沿x, 推导E\ensuremath{{}^{\prime\prime}}+k(z)\ensuremath{{}^{2}}E0; (c)\ensuremath{\omega}=\ensuremath{\omega}p在z0、n\ensuremath{\simeq}n0(1+z/a), 用Airy求通解; (d)低密侧WKB两支的E\ensuremath{-}/E+.
\end{exercise}
\begin{solution}
(a)匹配切向E/H, 非磁介质阻抗 $Z=Z_0/n_r$, 得\ensuremath{\Gamma}=(1\ensuremath{-}nr)/(1+nr), $n_r=\sqrt{1-\omega_p^2/\omega^2}$取衰减支. \ensuremath{\omega}<\ensuremath{\omega}p时nr纯正虚, $|\Gamma|=1$; n给fpe=8.9787MHz. 未给薄层厚度, 按单界面半无限等离子体处理, 若有限slab需第二界面.
(b)由Faraday与Amp\`ere消H, $E''+k_0^2[1-\omega_p^2(z)/\omega^2]E=0$. (c)\ensuremath{\zeta}=(k0\ensuremath{{}^{2}}/a)$^{1/3}$z, $E=C\operatorname{Ai}(\zeta)+D\operatorname{Bi}(\zeta)$, 高密侧不增长要求$D=0$. (d)低侧S(z)=$\int_z^0k\dd z$, Ai\ensuremath{\propto}sin(S+\ensuremath{\pi}/4). 取入射E+e$^{-iS}$/√k、反射E\ensuremath{-}e$^{iS}$/√k, 得\ensuremath{\Gamma}=\ensuremath{-}i、$|\Gamma|=1$. 改积分相位零点会改变复相位, 全反射能量不变; 不把sin渐近里的两项都称入射.
\end{solution}
